% !TEX program = xelatex
% 数学物理方法1 · 第四章专题 · 复数方程的根与极点判定
\documentclass[11pt]{ctexart}
%=====================================================================
% theme-violet.tex —— 讲义精讲 · 紫花主题
% 用法：\documentclass[11pt]{ctexart}
%       \input{assets/theme-violet}
%       \begin{document} ... \end{document}
% 只用 ctexart 自带能力 + xcolor/enumitem/fancyhdr/pifont/amsmath，
% 不依赖 tikz/tcolorbox，保证 MiKTeX 可直接编译。
%=====================================================================

\usepackage[a4paper,left=2.0cm,right=2.0cm,top=2.1cm,bottom=2.3cm,%
            headsep=12pt,footskip=22pt]{geometry}
\usepackage{amsmath,amssymb}
\usepackage{needspace}
\usepackage{xcolor}
\usepackage{array}
\usepackage{enumitem}
\usepackage{pifont}
\usepackage{fancyhdr}
\usepackage[hidelinks]{hyperref}

%---------------------------------------------------------------- 配色
\definecolor{violet}{RGB}{107,78,143}      % 主色：紫
\definecolor{violetdeep}{RGB}{63,44,92}    % 深紫：标题/强调
\definecolor{petal}{RGB}{183,159,208}      % 花瓣：装饰线、花饰
\definecolor{bandviolet}{RGB}{247,243,251} % 浅紫底：档案/提示框
\definecolor{bandcream}{RGB}{252,249,242}  % 米白底：易错/注意框
\definecolor{rulepurple}{RGB}{217,207,230} % 分隔线
\definecolor{inkgray}{RGB}{92,96,110}      % 次要文字

%---------------------------------------------------------------- 字体
% 有 Noto CJK 用 Noto，没有就退回 ctex 默认（Windows 下为宋体/黑体）
\IfFontExistsTF{Noto Serif CJK SC}{\setCJKmainfont{Noto Serif CJK SC}}{}
\IfFontExistsTF{Noto Sans CJK SC}{\setCJKsansfont{Noto Sans CJK SC}}{}

%---------------------------------------------------------------- 版式
\setlist{nosep,leftmargin=1.9em}
\linespread{1.28}
\setlength{\parskip}{3.5pt}
\setlength{\parindent}{0pt}
\setlength{\tabcolsep}{5pt}
\renewcommand{\arraystretch}{1.28}

%---------------------------------------------------------------- 花饰
% \ding{95} = 五瓣小花，\ding{96} = 星形花；都来自 pifont，必定存在
\newcommand{\orn}{\textcolor{petal}{\ding{95}}}          % 行内小标
\newcommand{\ornstar}{\textcolor{petal}{\ding{96}}}      % 星形花
\newcommand{\flowerrule}{%                                % 居中花饰分隔线
  \par\vspace{7pt}%
  \noindent\makebox[\linewidth]{\color{petal}%
    \rule{0.20\linewidth}{0.6pt}\hspace{8pt}{\small\ding{95}}\hspace{4pt}%
    \raisebox{1.1pt}{\tiny\ding{95}}\hspace{4pt}{\small\ding{95}}\hspace{8pt}%
    \rule{0.20\linewidth}{0.6pt}}%
  \par\vspace{7pt}}

%---------------------------------------------------------------- 标题与节头
\newcommand{\doctitle}[3]{%   #1 主标题  #2 副标题  #3 元信息（可用 \\ 分行）
  \begin{center}
    {\color{petal}\small \orn\enspace #2\enspace \ornstar}\\[7pt]
    {\LARGE\bfseries\color{violetdeep} #1}\\[10pt]
    {\color{petal}\rule{0.55\linewidth}{1pt}}\\[4pt]
    {\small\color{inkgray} #3}
  \end{center}\vspace{2pt}}

\newcommand{\hdr}[2]{%   #1 章节号（如 §0）  #2 标题
  \par\vspace{12pt}%
  \noindent\colorbox{violetdeep}{\color{white}\bfseries\small\;#1\;}\hspace{7pt}%
  {\color{violetdeep}\large\bfseries #2}\par\nobreak
  \vspace{3pt}{\color{petal}\hrule height 1.2pt}\vspace{6pt}}

\newcommand{\subhdr}[1]{%
  \par\vspace{6pt}%
  \noindent{\color{violet}\bfseries #1}\par\vspace{1pt}}

%---------------------------------------------------------------- 提示框
\newcommand{\callout}[1]{% 浅紫底：讲义档案、读法说明
  \par\medskip\noindent
  \fcolorbox{rulepurple}{bandviolet}{\parbox{\dimexpr\linewidth-2\fboxsep-2\fboxrule}{\vspace{2pt}#1\vspace{2pt}}}%
  \par\medskip}

\newcommand{\warnbox}[1]{% 米白底：易错、边界条件
  \par\medskip\noindent
  \fcolorbox{rulepurple}{bandcream}{\parbox{\dimexpr\linewidth-2\fboxsep-2\fboxrule}{\vspace{2pt}#1\vspace{2pt}}}%
  \par\medskip}

\newcommand{\gapitem}[2]{% #1 缺口类型（紫） #2 内容
  \par\vspace{2pt}\noindent
  \hangindent=1.4em\hangafter=1
  {\color{petal}\small\ding{96}}\enspace{\bfseries\color{violetdeep}#1}\enspace #2\par}

% 题干块（\qstem）：每道例题／作业的「详解」开头必须先复述题干与出处。
% 起因：2026-10-25 第五章讲义把七道作业的题面集中放在一页、详解另起一页，
% 读者翻到「第 4 题」时看不到条件，等于一道没有题干的题目。
% 规则：详解块内自带题干，读者不翻页就能读懂题在问什么。
\newcommand{\qstem}[1]{%
  \par\medskip\noindent
  \fcolorbox{violet}{bandviolet}{%
    \parbox{\dimexpr\linewidth-2\fboxsep-2\fboxrule}{%
      \vspace{2pt}%
      {\small\bfseries\color{violet}题干}\hspace{0.8em}#1%
      \vspace{2pt}}}%
  \par\medskip}

% v2: 保护题干与开头，长题可跨页，续页补上下文。
\newcommand{\qbreak}{\needspace{6\baselineskip}}
% 独立练习：ID、来源、题干、考点、答案标签（通常 ans:ID）。
\newcommand{\practice}[5]{%
  \qbreak\subhdr{练习 #1}\qstem{（#2）#3}%
  \textbf{考点}：#4\par
  详解见附录第~\pageref{#5}~页。\par\medskip}
% 附录开头：ID、来源、条件摘要；自动产生 ans:ID 标签。
\newcommand{\answerstart}[3]{%
  \qbreak\subhdr{#1\quad 详解}\label{ans:#1}\qstem{（#2）#3}}
\newcommand{\answercontinue}[2]{%
  \needspace{5\baselineskip}\subhdr{#1\quad 续解}%
  {\small\color{inkgray}本页接续：#2}\par\medskip}
% 延迟提供默认值，让旧讲义可在导言区保留自己的 \newcommand{\step}。
\AtBeginDocument{\providecommand{\step}[1]{\needspace{3\baselineskip}\par\medskip\textbf{#1}\par}}

% 前提旁注：只放在「需要理解的那一步」后面，给陈述 + 本题代入，不重证。
\newcommand{\aside}[1]{%
  \par\vspace{1pt}\noindent
  \hangindent=1.6em\hangafter=1
  {\small\color{inkgray}%
    \hspace*{0.6em}{\color{petal}$\triangleright$}\hspace{2pt}#1}%
  \par\vspace{2pt}}

%---------------------------------------------------------------- 表格
% 本机 MiKTeX 的 colortbl 与 array 版本不匹配（\rowcolor 会报错），
% 故表格统一用「上粗紫线 + 表头 + 中细线」的三线式，不依赖 colortbl。
\newcommand{\tabtop}{\par\vspace{3pt}{\color{violet}\hrule height 1.2pt}\vspace{1pt}}
\newcommand{\tabmid}{\par\vspace{1pt}{\color{petal}\hrule height 0.5pt}\vspace{1pt}}
\newcommand{\tabbot}{\par\vspace{3pt}{\color{rulepurple}\hrule height 0.5pt}\vspace{8pt}}
\newcommand{\thead}[1]{{\bfseries\color{violetdeep}#1}}

%---------------------------------------------------------------- 页眉页脚
\pagestyle{fancy}
\fancyhf{}
\renewcommand{\headrulewidth}{0.4pt}
\renewcommand{\headrule}{\hbox to\headwidth{\color{rulepurple}\leaders\hrule height \headrulewidth\hfill}}
\fancyhead[L]{\small\color{inkgray}\songti \doctag}
\fancyhead[R]{\small\color{inkgray}\songti \coursetag}
\fancyfoot[C]{\small\color{inkgray}\orn\enspace{\color{violet}\thepage}\enspace\orn}

% 每篇讲义在导言区设这两个值
\newcommand{\doctag}{}
\newcommand{\coursetag}{}
\newcommand{\settags}[2]{\renewcommand{\doctag}{#1}\renewcommand{\coursetag}{#2}}

%---------------------------------------------------------------- 数学记号
\newcommand{\vect}[1]{\boldsymbol{#1}}
\newcommand{\dd}{\,\mathrm{d}}
\newcommand{\ddx}{\frac{\mathrm{d}}{\mathrm{d}x}}
\addtolength{\jot}{2pt}



\usepackage{tikz}
\usetikzlibrary{arrows.meta,positioning}
\settags{第四章专题 · 复数方程的根与极点判定}{数学物理方法1}
% p{} 列首的 \color 会触发换段，改用 \textcolor
\renewcommand{\thead}[1]{\textbf{\textcolor{violetdeep}{#1}}}

\renewcommand{\qbreak}{\needspace{0.32\textheight}}
\let\origsubhdr\subhdr
\renewcommand{\subhdr}[1]{\par\needspace{4\baselineskip}\origsubhdr{#1}}
\let\orighdr\hdr
\renewcommand{\hdr}[2]{\par\orighdr{#1}{#2}}
\setlength{\abovedisplayskip}{7pt plus 2pt minus 2pt}
\setlength{\belowdisplayskip}{7pt plus 2pt minus 2pt}

\begin{document}

\doctitle{第四章专题　复数方程的根与极点判定}
{从 $z^{4}-1=0$ 到“几阶极点、在不在圈内”}
{依据《第四章 留数定理及应用》上传课件（78 页）p11--p12、p18--p23、p51\quad 2026-11-02\quad 数学物理方法1}

\hdr{§0}{缘起：一行字里压着三件事}

课件 p18--p19 的例 2 只有一句判据：“函数 $f(z)=\dfrac{z}{z^{4}-1}$ 有四个一阶极点 $z=\pm1,\ \pm i$，且在积分半径内，则规则 II”。这句话读起来像结论，其实是三步操作叠在一起的结果：解方程 $z^{4}-1=0$ 拿到四个根；看每个根是几重；判断四个根在 $|z|=2$ 里面还是外面。三步各自都不难，可它们被压在同一个句子里，任何一步没做熟，整道题就停在第一行。

后面几页是同一个结构。p20 写“$z=\pm i$ 是该函数的两个二阶极点”，要的是看出 $(z^{2}+1)^{2}=(z-i)^{2}(z+i)^{2}$；p21 写“可以看出，该函数有两个极点，分别为三阶极点 $z=0$ 和一阶极点 $z=i$”，要的是读出分母各因式的幂次；p51 写“$z^{4}+a^{4}=0$，即 $z^{4}=-a^{4}=a^{4}e^{i(2k+1)\pi}$”，这一处反而把求根写全了。同一章里，p51 写得完整，p18--p19 直接给结果——两处其实是同一套方法，只是讲法一处细、一处粗。

这份讲义把那套方法拆成四步，每一步配一个能自己复核的检查。目标只有一个：拿到一个有理函数，能在三十秒内写出“极点清单”——每个极点的位置、阶数、在不在给定围道内——并且知道这份清单为什么可信。

\subhdr{0.1\quad 记号（全文统一）}

\begin{itemize}
\item $P,\ Q$：多项式；$P$ 在分子，$Q$ 在分母，$f=P/Q$。
\item $z_0$：待判断的点；$m$：根（或极点）的重数（阶数）；$n$：方程的次数。
\item $\arg z$：辐角，多值；$\operatorname{Arg}z\in(-\pi,\pi]$：主值。求根时用哪个都可以，见 §4.3。
\item $\operatorname{Res}f(z_0)$：留数；$C$：正向简单闭曲线；$R$：围道 $|z|=R$ 的半径。
\end{itemize}

\subhdr{0.2\quad 五处要补的缺口（逐条见 §13）}

\gapitem{代数缺口 · p19}{四个极点直接列出，没有解 $z^{4}-1=0$ 的过程。补在 §4、§5.3。}

\gapitem{重数缺口 · p20}{由 $(z^{2}+1)^{2}$ 直接跳到“二阶极点”，没有说“幂次就是阶数”。补在 §7.1、§7.2。}

\gapitem{压缩缺口 · p21}{“可以看出”三个字吞掉了“分母因式的指数相加”。补在 §7.2。}

\gapitem{判圈缺口 · p22--p23}{两行不等式给出 $|z_1|>1>|z_2|$，没有出处。补在 §2.2、§7.4。}

\gapitem{结构缺口 · p51}{四个根写得完整，但没有指出这就是 $z^{n}=a$ 的通式。补在 §4。}

\hdr{§1}{三个必须先立住的事实}

这一节只有三条，但它们是后面所有操作的合法性来源：第一条告诉你“一共几个根”，第二条告诉你“一个根算几重”，第三条告诉你“实系数方程可以少解一半”。

\subhdr{1.1\quad 代数基本定理与因式分解形式}

\textbf{定理（代数基本定理）}：设 $n\ge1$，$a_n\ne0$，$P(z)=a_nz^{n}+a_{n-1}z^{n-1}+\cdots+a_1z+a_0$ 为复系数多项式，则在复平面 $\mathbb{C}$ 中至少存在一点 $z_0$，使 $P(z_0)=0$。

\textbf{推论（因式分解形式）}：把定理反复用在降次后的多项式和商上，得
\begin{equation}
P(z)=a_n(z-z_1)^{m_1}(z-z_2)^{m_2}\cdots(z-z_s)^{m_s},
\qquad m_1+m_2+\cdots+m_s=n,
\label{eq:factor}
\end{equation}
其中 $z_1,\dots,z_s$ 互不相同，各 $m_j$ 为正整数。

它的实用形式只有一句：\textbf{$n$ 次方程一共有 $n$ 个根，重根按重数算}。$(z^{2}+1)^{2}=0$ 的次数是 4，不同的根只有 $\pm i$ 两个，但 $i$ 算两次、$-i$ 算两次，合起来正好 4 个。做题时用这条检查：根解完数一遍，数目对不上，就是漏了根或者漏了重数。

\aside{定理只保证根存在，不告诉你根在哪、怎么求。所以它给的是记账规则（一共几个、重数之和等于次数），具体求根仍要按 §2 到 §6 的办法做。}

\subhdr{1.2\quad 重数的三个等价说法}

“$z_0$ 是 $P$ 的 $m$ 重根”与下面两条等价：

(i) $P(z_0)=P'(z_0)=\cdots=P^{(m-1)}(z_0)=0$，而 $P^{(m)}(z_0)\ne0$；

(ii) $P(z)=(z-z_0)^{m}h(z)$，其中 $h$ 在 $z_0$ 解析且 $h(z_0)\ne0$。

(ii) 是“因式分解”的严格版本，(i) 是它的计算版本。最常用的只是 (i) 的一阶特例：\textbf{$P(z_0)=0$ 且 $P'(z_0)\ne0$，就是单根}；若 $P'(z_0)$ 也为零，就至少是二重。这一条在后面判极点阶数时会反复用到。

\subhdr{1.3\quad 实系数的两个对称性}

若 $P$ 的系数全是实数，则对任意 $z$ 有 $P(\bar z)=\overline{P(z)}$。由此得两件事：非实根必成共轭对出现，且重数相同；奇次实系数多项式至少有一个实根（复根成对消耗偶数次，剩下的奇数次只能由实根填满）。

\aside{这两条不是装饰：它们让“解出一半根”成为合法操作。四次实系数方程只要算出两个根，另外两个直接取共轭，不必再解。}

\hdr{§2}{二次方程：公式照用，开方要开对}

二次式在整条链条里出现得最多，因为降次到最后总剩一个二次式。它的公式在复数里与实数完全一样，只有“开方”这一处需要当心。

\subhdr{2.1\quad 求根公式与复数的开方}

设 $a\ne0$，$b,c\in\mathbb{C}$，$az^{2}+bz+c=0$，则
\begin{equation}
z_{1,2}=\frac{-b\pm\sqrt{b^{2}-4ac}}{2a}.
\label{eq:quad}
\end{equation}

在实数里 $\sqrt{\Delta}$ 约定取非负值，$\pm$ 给出两个根；在复数里 $\sqrt{\Delta}$ 本身就有两个值，与 $\pm$ 重复。正确的用法是先把判别式写成极坐标，再开方：

设 $\Delta=b^{2}-4ac=|\Delta|e^{i\varphi}$，取 $\sqrt{\Delta}=|\Delta|^{1/2}e^{i\varphi/2}$（另一个值是它的相反数，已由 $\pm$ 覆盖），代回式 \eqref{eq:quad}。

\textbf{例}：$z^{2}+2z+2=0$。$\Delta=4-8=-4=4e^{i\pi}$，故 $\sqrt{\Delta}=2e^{i\pi/2}=2i$，于是
\[ z=\frac{-2\pm2i}{2}=-1\pm i . \]
核对：$z=-1+i$ 时 $z^{2}=-2i$，$2z=-2+2i$，与常数 $2$ 相加正好为 $0$。

\aside{为什么不能直接写 $\sqrt{-4}=2i$ 就完事：这只是在“取了某一个平方根”，而式 \eqref{eq:quad} 里的 $\pm$ 会把两个平方根都用到。把 $\Delta$ 写成极坐标，是为了让“取哪个平方根”变成看得见的 $\varphi/2$，不再靠猜。}

\subhdr{2.2\quad 判别式的三种情形（实系数）}

$\Delta>0$：两个不同实根。$\Delta=0$：一个二重实根 $z=-b/(2a)$。

$\Delta<0$：一对共轭复根。这时由韦达定理 $z_1z_2=c/a>0$ 且 $z_2=\bar z_1$，所以
\[ |z_1|^{2}=z_1\bar z_1=z_1z_2=\frac{c}{a},
\qquad |z_1|=|z_2|=\sqrt{c/a}. \]
这个式子直接回答“这对根在不在半径 $R$ 的圆内”，不必解方程。

\textbf{例（课件 p22--p23 例 5 的判圈步骤）}：$\varepsilon z^{2}+2z+\varepsilon=0$，$0<\varepsilon<1$。这里 $\Delta=4-4\varepsilon^{2}>0$，所以两根都是实数；韦达定理给
\[ z_1z_2=1,\qquad z_1+z_2=-\frac{2}{\varepsilon}<-2 . \]
两个负数乘积为 $1$、和为 $-2/\varepsilon$：设 $z_1=-s$、$z_2=-1/s$（$s>0$），则 $z_1+z_2=-(s+1/s)<-2$，说明 $s\ne1$，于是一个模大于 $1$、一个模小于 $1$。课件只写了结论 $|z_1|>1>|z_2|$，这两行就是它的依据。

\aside{若 $\Delta<0$ 而系数是复数，则 $z_2\ne\bar z_1$，$|z|^{2}=c/a$ 不再成立，只能老老实实算出两个根再取模。}

\hdr{§3}{三次以上：没有公式可背，只有三条路}

三次、四次方程有根式解（Cardano、Ferrari 公式），但计算繁，本章用不上；五次以上没有根式解（Abel--Ruffini 定理）。所以做留数题时，遇到三次以上的分母，实际只有三条路：

\begin{itemize}
\item \textbf{路线 A · 结构识别}：式子是平方差、立方差、双二次或某种完全平方，一比就分解开（§6）。
\item \textbf{路线 B · 二项式型}：方程形如 $z^{n}=a$。这是唯一有统一公式的一类，也是本章出现频率最高的一类（§4）。
\item \textbf{路线 C · 猜一个根再降次}：用有理根定理筛出一个小整数根，除掉一个一次因子，次数降一，直到剩下二次（§5）。
\end{itemize}

三条路不是互斥的：$z^{4}-1=0$ 三条都能走，结果必须一致——这种“两条路给同一组根”的一致性，就是最好的自检。

做题时的顺序建议：\textbf{先看是不是 $z^{n}=a$}（最快，四五行出全部根），不是再看有没有明显的平方差或完全平方，最后才试根。

\aside{为什么把 $z^{n}=a$ 排第一：它的全部根一次性给出，而且所有根的模相同，判“在不在圈内”只需比一个数。本章的 $z^{4}-1$、$z^{4}+a^{4}$、$z^{3}=8$ 都属于这一类。}

\hdr{§4}{路线 B：$z^{n}=a$ 的全部根}

\subhdr{4.1\quad 公式与推导}

设 $a\ne0$，写成极坐标 $a=|a|e^{i\theta}$，其中 $\theta$ 是 $a$ 的某一个辐角。设 $z=re^{i\varphi}$（$r>0$，$\varphi$ 为实数），则 $z^{n}=r^{n}e^{in\varphi}$。两个复数相等，必须模相等、辐角相差 $2\pi$ 的整数倍：

\[ r^{n}=|a|,\qquad n\varphi=\theta+2k\pi\ (k\in\mathbb{Z}). \]

第一式给 $r=|a|^{1/n}$（正实数，唯一），第二式给 $\varphi=\dfrac{\theta+2k\pi}{n}$。当 $k$ 取 $0,1,\dots,n-1$ 时，这 $n$ 个角互不相同；再大的 $k$ 只是重复已有的角。于是
\begin{equation}
z_k=|a|^{1/n}\exp\left(i\,\frac{\theta+2k\pi}{n}\right),
\qquad k=0,1,\dots,n-1 .
\label{eq:roots}
\end{equation}

两条立刻可用的结论：

(i) \textbf{所有根的模都等于 $|a|^{1/n}$}，与 $k$ 无关；

(ii) 相邻两根的辐角差都等于 $2\pi/n$。

把两条合起来：$z^{n}=a$ 的 $n$ 个根，是内接于圆 $|z|=|a|^{1/n}$ 的正 $n$ 边形，第一个顶点在辐角 $\theta/n$ 处。

\begin{center}
\begin{tikzpicture}[scale=1.05]
\draw[->,gray] (-2.7,0) -- (2.7,0) node[right,black]{$\mathrm{Re}\,z$};
\draw[->,gray] (0,-2.3) -- (0,2.3) node[above,black]{$\mathrm{Im}\,z$};
\draw[violet,thick] (0,0) circle (2);
\foreach \k in {0,1,2,3,4}{
  \coordinate (P\k) at ({2*cos(30+72*\k)},{2*sin(30+72*\k)});
  \fill[violet] (P\k) circle (2.2pt);
}
\draw[gray,dashed] (0,0) -- (P0);
\draw[gray,dashed] (0,0) -- (P1);
\draw[->,thick,violetdeep] (0.85,0) arc (0:30:0.85);
\node at (1.30,0.16) {$\theta/n$};
\node[left] at (-0.30,1.03) {$\small r$};
\node[below right] at (P0) {\small $z_0$};
\node[above left] at (P1) {\small $z_1$};
\end{tikzpicture}
\end{center}

\textbf{图 1 读法}：圆心在原点，半径 $r=|a|^{1/n}$ 由 $a$ 的模开 $n$ 次方得到；第一个根 $z_0$ 与正实轴的夹角是 $\theta/n$（不是 $\theta$）；从 $z_0$ 起每次转 $2\pi/n$ 得到下一个。图里画的是 $n=5$：虚线是 $z_0$ 与 $z_1$ 的半径，两线夹角 $72^\circ$；从正实轴到 $z_0$ 的小弧是 $\theta/n$。$n=4$ 时这个正五边形变成正方形（下一小节的表第一行），$n=2$ 时退化成一对相反点——那就是二次方程的两个根。

\subhdr{4.2\quad $\theta$ 怎么取（这一步最容易错）}

$\theta$ 只要\textbf{是 $a$ 的一个真辐角}就行，取 $\theta+2\pi$ 也照样对，因为 $+2k\pi$ 已经把多值吸收了。错的是“就近取 $0$”：

\begin{itemize}
\item $a>0$ 实数：$\theta=0$；
\item $a<0$ 实数：$\theta=\pi$，\textbf{不是 $0$}；
\item $a=it$（$t>0$）：$\theta=\pi/2$；$a=-it$：$\theta=-\pi/2$（取 $3\pi/2$ 得到同一组根）；
\item 一般 $a=x+iy$：取 $\theta=\operatorname{Arg}a$，即 $a$ 所在象限的那个角。
\end{itemize}

最常见的错误是把 $z^{4}=-1$ 的 $\theta$ 写成 $0$，于是解出 $z=\pm1,\pm i$——那其实是 $z^{4}=1$ 的根。差一个 $\pi$，整套答案就换了另一道题。

\subhdr{4.3\quad 五个必须一眼看出的方程}

这张表只是把式 \eqref{eq:roots} 在五个方程上各读一遍：先定 $|a|$ 与 $\theta$，再套模与辐角。

\begin{minipage}{\linewidth}
\small
\tabtop
\begin{tabular}{@{}p{0.15\linewidth}p{0.17\linewidth}p{0.44\linewidth}p{0.12\linewidth}@{}}
\thead{方程} & \thead{根的模} & \thead{根} & \thead{邻根角差} \\
\end{tabular}
\tabmid
\begin{tabular}{@{}p{0.15\linewidth}p{0.17\linewidth}p{0.44\linewidth}p{0.12\linewidth}@{}}
$z^{4}=1$ & $1$ & $1,\ i,\ -1,\ -i$ & $90^\circ$ \\
$z^{4}=-1$ & $1$ & $e^{i\pi/4},\ e^{3i\pi/4},\ e^{5i\pi/4},\ e^{7i\pi/4}$ & $90^\circ$ \\
$z^{3}=8$ & $2$ & $2,\ -1+i\sqrt3,\ -1-i\sqrt3$ & $120^\circ$ \\
$z^{4}=-16$ & $2$ & $2e^{i\pi/4},\ 2e^{3i\pi/4},\ 2e^{5i\pi/4},\ 2e^{7i\pi/4}$ & $90^\circ$ \\
$z^{5}=1+i$ & $2^{1/10}\approx1.072$ & $2^{1/10}e^{i(\pi/20+2k\pi/5)}$，$k=0,\dots,4$ & $72^\circ$ \\
\end{tabular}
\tabbot
\end{minipage}

\textbf{逐行读法}：第一行 $a=1=1\cdot e^{i0}$，模 $1^{1/4}=1$，角 $0+\frac{2k\pi}{4}$，得 $0^\circ,90^\circ,180^\circ,270^\circ$。第二行 $a=-1=e^{i\pi}$，模仍是 $1$，角 $\frac{\pi+2k\pi}{4}=\frac{\pi}{4}+\frac{k\pi}{2}$，所以四个点整体转了 $45^\circ$——这是 $z^{4}=1$ 与 $z^{4}=-1$ 的全部差别。第三行 $a=8$，模 $8^{1/3}=2$，角 $0,\frac{2\pi}{3},\frac{4\pi}{3}$；第二个根写成 $-1+i\sqrt3$ 后模仍是 $2$（$\sqrt{1+3}=2$），可以当场验算。第四行 $a=-16=16e^{i\pi}$，模 $16^{1/4}=2$，角 $\frac{\pi}{4}+\frac{k\pi}{2}$；第一个根 $2e^{i\pi/4}=\sqrt2(1+i)$。第五行 $a=1+i=\sqrt2\,e^{i\pi/4}$，模 $2^{1/10}\approx1.072$，角 $\frac{\pi}{20}+\frac{2k\pi}{5}$，第一个根离正实轴只有 $9^\circ$。

\textbf{边界行}：$a=0$ 时方程退化为 $z^{n}=0$，只有一个根 $z=0$，但按重数是 $n$ 重（$z^{n}=(z-0)^{n}$），不是 $n$ 个不同根——这一行的模写成 $0$，代入式 \eqref{eq:roots} 会得到同一点重复 $n$ 次，与式 \eqref{eq:factor} 的记账规则一致。$a$ 为负实数时 $\theta=\pi$，第一个顶点的角是 $\pi/n$ 而不是 $0$。

\subhdr{4.4\quad 单位根与 $z^{n}-1$ 的分解}

$z^{n}=1$ 的 $n$ 个根叫 $n$ 次单位根，记 $\omega_k=e^{2k\pi i/n}$。设 $\omega\ne1$ 是一个 $n$ 次单位根，则 $\omega^{n}=1$，由等比数列求和
\[ 1+\omega+\omega^{2}+\cdots+\omega^{n-1}
=\frac{\omega^{n}-1}{\omega-1}=0 . \]
把 $\omega$ 换回变量 $z$，同一段代数就是因式分解：
\begin{equation}
z^{n}-1=(z-1)\left(z^{n-1}+z^{n-2}+\cdots+z+1\right).
\label{eq:unit}
\end{equation}
式 \eqref{eq:unit} 是路线 C 的常用工具：$z^{4}-1=(z-1)(z^{3}+z^{2}+z+1)$ 正是它的 $n=4$ 特例。

\hdr{§5}{路线 C：猜一个根，降一次次}

\subhdr{5.1\quad 有理根定理：把“猜”变成“筛”}

\textbf{定理（有理根定理）}：设 $P(z)=a_nz^{n}+a_{n-1}z^{n-1}+\cdots+a_1z+a_0$ 是整系数多项式（$a_n\ne0$）。若有理数 $\dfrac{p}{q}$（$p,q$ 互素）是 $P$ 的根，则 $p\mid a_0$，$q\mid a_n$。

首项系数为 $1$ 时，$q\mid1$ 逼出 $q=1$，于是\textbf{首一整系数多项式的有理根必是整数，且必整除常数项}。这就是试根的合法名单：不是乱猜，是从常数项的因数里挑。

\aside{推导两行：把 $z=p/q$ 代入，两边乘 $q^{n}$，得 $a_np^{n}+a_{n-1}p^{n-1}q+\cdots+a_1pq^{n-1}+a_0q^{n}=0$。除最后一项外每项都含因子 $p$，故 $p\mid a_0q^{n}$，由 $p,q$ 互素得 $p\mid a_0$；同理把式子看成“除第一项外每项含 $q$”，得 $q\mid a_n$。}

\subhdr{5.2\quad 降次：除一个一次因子}

试出一个根 $z_1$ 后，用 $(z-z_1)$ 除 $P$（长除法或综合除法），因为整除，余式为 $0$，商是 $n-1$ 次多项式。次数降一，继续做。若商在 $z_1$ 处仍然为 $0$，说明 $z_1$ 是重根，要再除一次——这一步就是重数容易丢的地方。

示范（$z^{4}-1$ 除以 $z-1$）：被除式系数 $1,0,0,0,-1$，综合除法依次得商系数 $1,1,1,1$、余 $0$，即
\[ z^{4}-1=(z-1)\left(z^{3}+z^{2}+z+1\right). \]

\subhdr{5.3\quad 完整示范：$z^{4}-1=0$}

第一步，试根名单：$a_0=-1$ 的因数是 $\pm1$，只试这两个。

$z=1$：$1-1=0$，是根，除得 $(z-1)(z^{3}+z^{2}+z+1)$。

$z=-1$：$-1+1-1+1=0$，是根，再除 $(z+1)$，得 $(z-1)(z+1)(z^{2}+1)$。

第二步，处理剩下的二次式：$z^{2}+1=0$ 即 $z^{2}=-1$，由 §2 得 $z=\pm i$。

于是 $z^{4}-1=(z-1)(z+1)(z-i)(z+i)$，四个根 $1,-1,i,-i$ 互不相同，全是单根。

\aside{用 §4 核一遍：$z^{4}=1$ 由式 \eqref{eq:roots} 一步给出 $1,i,-1,-i$，与上面一致。两条路都要会：路线 B 快，路线 C 通用——遇到 $z^{3}-3z+2$ 这种不是二项式的方程，只有路线 C 走得通。}

\subhdr{5.4\quad 重根示范：$z^{3}-3z+2=0$}

试根名单：$a_0=2$ 的因数是 $\pm1,\pm2$。

$z=1$：$1-3+2=0$，是根，除 $(z-1)$ 得 $z^{2}+z-2$。

$z^{2}+z-2=(z+2)(z-1)$，所以
\[ z^{3}-3z+2=(z-1)^{2}(z+2), \]
$z=1$ 是\textbf{二重根}。用 §1.2 的导数检验复核：$P'(z)=3z^{2}-3$，$P'(1)=0$；$P''(z)=6z$，$P''(1)=6\ne0$，故恰好二重。

\warnbox{本章最容易丢分的一步就在这里：$z=1$ 试出来就停手，会把它当单根，极点阶数少记一级，规则 III 也用不上，最后差一个系数。降次后若商在同一点仍为 $0$，必须再除一次。}

\hdr{§6}{路线 A：结构识别}

有些分母不需要“解”，看结构就能拆开。四种最常用。

\subhdr{6.1\quad 平方差与立方差}

$z^{4}-1=(z^{2}-1)(z^{2}+1)=(z-1)(z+1)(z-i)(z+i)$。

$z^{6}-1=(z^{3}-1)(z^{3}+1)$，两个立方差再各拆一次。

\subhdr{6.2\quad 双二次：令 $u=z^{2}$}

形如 $z^{4}+pz^{2}+q=0$ 的方程，先解 $u^{2}+pu+q=0$ 得 $u_1,u_2$，再对每个 $u$ 解 $z^{2}=u$；每个 $u$ 给两个根，合起来四个。

\textbf{例}：$z^{4}-5z^{2}+4=0$。$u^{2}-5u+4=0$ 给 $u=1,4$；$z^{2}=1$ 给 $\pm1$，$z^{2}=4$ 给 $\pm2$。四个根 $1,-1,2,-2$。

\subhdr{6.3\quad 完全平方：重根从哪里来}

$z^{4}+2z^{2}+1=(z^{2}+1)^{2}$。这一行就是课件 p20 里“二阶极点”的来源：分母是某个式子的平方，每个根就变成二重根。

同理 $z^{4}-2z^{2}+1=(z^{2}-1)^{2}$，$z^{6}+2z^{3}+1=(z^{3}+1)^{2}$。

\subhdr{6.4\quad 两个“看起来像 $z^{4}+1$”的式子}

$z^{4}+4$ 不能直接写成平方差，但可以配成完全平方再减：
\[ z^{4}+4=(z^{2}+2)^{2}-(2z)^{2}=(z^{2}-2z+2)(z^{2}+2z+2), \]
这叫 Sophie Germain 分解。两个二次式的根分别是 $1\pm i$ 与 $-1\pm i$，四个根的模都是 $\sqrt2$。

$z^{4}+1$ 的配法是 $(z^{2}+1)^{2}-2z^{2}$：
\[ z^{4}+1=(z^{2}-\sqrt2\,z+1)(z^{2}+\sqrt2\,z+1). \]
用式 \eqref{eq:quad} 算 $z^{2}-\sqrt2z+1=0$：
\[ z=\frac{\sqrt2\pm\sqrt{2-4}}{2}=\frac{\sqrt2\pm i\sqrt2}{2}=e^{\pm i\pi/4}. \]
与 §4 对 $z^{4}=-1$ 给出的 $e^{i\pi/4},e^{3i\pi/4},\dots$ 完全一致——两条路必须给同一组根，这是最好的自检。

\aside{两个式子的差别只在常数项：$z^{4}+4$ 的根在半径 $\sqrt2$ 的圆上，$z^{4}+1$ 的根在单位圆上。做题时先算模，一眼就能分开。}

\hdr{§7}{从根到极点：位置与阶数}

根解出来只是原料，极点清单还要两样信息：每个极点的阶数、它在不在围道内。

\subhdr{7.1\quad 阶数判据（完整形式）}

设 $f(z)=\dfrac{P(z)}{Q(z)}$，$P,Q$ 在 $z_0$ 的邻域内解析，$Q(z_0)=0$。设 $Q$ 在 $z_0$ 有 $m$ 阶零点，$P$ 在 $z_0$ 有 $k$ 阶零点（约定 $P(z_0)\ne0$ 时 $k=0$）。则
\begin{equation}
\begin{cases}
k\ge m: & z_0\ \text{是}\ f\ \text{的可去奇点，不是极点};\\[2pt]
k<m:  & z_0\ \text{是}\ f\ \text{的}\ m-k\ \text{阶极点}.
\end{cases}
\label{eq:order}
\end{equation}

最常用的是 $k=0$ 那一支：分子在 $z_0$ 不为零、分母有 $m$ 阶零点，则 $z_0$ 就是 $m$ 阶极点。\textbf{分母因式的指数就是阶数}——这句话在有重因子时最省事。

\subhdr{7.2\quad 三条实操检验（互相复核）}

\textbf{检验一 · 因式分解}（最稳）：把分母写成 $(z-z_0)^{m}$ 与其余因子之积，$m$ 即阶数。$Q=(z^{2}+1)^{2}=(z-i)^{2}(z+i)^{2}$，读指数直接得两个二阶极点。

\textbf{检验二 · 导数检验}：若 $Q(z_0)=Q'(z_0)=\cdots=Q^{(m-1)}(z_0)=0$ 而 $Q^{(m)}(z_0)\ne0$，则 $z_0$ 是 $Q$ 的 $m$ 阶零点。以 $Q=(z^{2}+1)^{2}$ 在 $z=i$ 为例：$Q(i)=0$；$Q'(z)=4z(z^{2}+1)$，$Q'(i)=0$；$Q''(z)=12z^{2}+4$，$Q''(i)=-8\ne0$，故二阶。

\textbf{检验三 · 极限实测}：若 $\displaystyle\lim_{z\to z_0}(z-z_0)^{k}f(z)$ 为有限非零数，则阶数恰为 $k$；若为 $0$，阶数小于 $k$；若为 $\infty$，阶数大于 $k$。不想分解时这条最快。

\subhdr{7.3\quad 一阶极点与规则 II（课件 p11--p12 给的就是这个形式）}

\textbf{定理（规则 II）}：设 $f=\dfrac{P}{Q}$，$P,Q$ 在 $z_0$ 解析，$P(z_0)\ne0$，$Q(z_0)=0$，$Q'(z_0)\ne0$，则 $z_0$ 是 $f$ 的一阶极点，且
\begin{equation}
\operatorname{Res}f(z_0)=\frac{P(z_0)}{Q'(z_0)} .
\label{eq:rule2}
\end{equation}

条件 $Q'(z_0)\ne0$ 正是“$z_0$ 是 $Q$ 的单根”的导数说法（§1.2）。所以“四个根互不相同”这句话，可以直接翻译成“四个一阶极点”，一步到位。

\aside{推导一行：$Q(z)=(z-z_0)\left[Q'(z_0)+O(z-z_0)\right]$，于是 $(z-z_0)f\to P(z_0)/Q'(z_0)$，再用一阶极点的定义 $\operatorname{Res}=\lim_{z\to z_0}(z-z_0)f$。}

\subhdr{7.4\quad 位置：圈内、圈外、还是圈上}

判据只有一句：算 $|z_k|$，与 $R$ 比。$|z_k|<R$ 在内部，要算进留数和；$|z_k|>R$ 在外部，一个都不算；$|z_k|=R$ 落在围道上，留数定理不能直接用（7.5）。

四种常见情形：

(a) $z^{n}=a$：全部根同模 $|a|^{1/n}$，一次判断定全体。$z^{4}=1$ 的四根模都是 $1$，所以对 $|z|=2$ 全在内部，对 $|z|=1$ 全在围道上——两个结论只差一次比较。

(b) 实系数二次、$\Delta<0$：共轭对，$|z|^{2}=c/a$（§2.2）。

(c) 已分解成一次因子：直接读出模，例如 $(z-1)(z+1)(z-2)(z+2)$ 的模是 $1,1,2,2$。

(d) 实系数二次、$\Delta>0$：两根实数，用韦达定理的乘积与和定性判断（§2.2 的例）。

\subhdr{7.5\quad 根正好落在围道上}

这时留数定理的前提（奇点严格在内部）不满足，硬套会得到一个错误的值。两种处置：把围道改小或改大，使根不再穿过；或者保留围道，在根处挖一个小半圆绕过去，此时贡献是 $i\pi\operatorname{Res}$ 而不是 $2\pi i\operatorname{Res}$（半留数，见第四章 p63--p70 补充内容）。

课件 p18--p19 的例 2 之所以取 $|z|=2$ 而不是 $|z|=1$，原因正在这里：四个根全落在单位圆上。

\textbf{一个小例子}：$\displaystyle\oint_{|z|=1}\frac{\dd z}{z^{2}-1}$。两个极点 $z=\pm1$ 都落在围道上（模都是 $1$）。不能直接套留数定理；若改取 $|z|=R$，则 $0<R<1$ 时圈内没有极点、积分为 $0$，$R>1$ 时两个极点都在内部、留数之和为 $0$、积分也是 $0$。坚持在原围道上算，就在 $\pm1$ 处各挖一个小半圆取主值，两个半留数 $\pm\frac{i\pi}{2}$ 正好抵消，结果仍是 $0$。

\needspace{0.34\textheight}
\hdr{§8}{主例题全程（课件 p18--p19 例 2）}

\qbreak
\qstem{（课件 p18--p19 例 2）计算 $\displaystyle I=\oint_{|z|=2}\frac{z}{z^{4}-1}\dd z$，要求分别用柯西公式与留数定理计算。}

\textbf{第一步 · 解方程}。$z^{4}-1=0$ 即 $z^{4}=1$。$a=1=1\cdot e^{i0}$，由式 \eqref{eq:roots}：模 $1^{1/4}=1$，角 $\frac{0+2k\pi}{4}$，得
\[ z=1,\ i,\ -1,\ -i . \]
也可用 §5.3 的路线 C 得到同一组根。

\textbf{第二步 · 数重数}。四个根互不相同；$Q'(z)=4z^{3}$ 在每个根处等于 $4z_k^{3}\ne0$，所以每个都是单根，对应一阶极点。这就是课件那句“四个一阶极点”的来历。

\textbf{第三步 · 判圈}。四个根的模都是 $1$，而 $1<2$，所以四个极点全在 $|z|=2$ 内部。

\begin{center}
\begin{tikzpicture}[scale=0.85]
\draw[->,gray] (-3.3,0) -- (3.3,0) node[right,black]{$\mathrm{Re}\,z$};
\draw[->,gray] (0,-2.15) -- (0,3.0) node[above,black]{$\mathrm{Im}\,z$};
\draw[gray,thick] (0,0) circle (1);
\draw[violet,thick] (0,0) circle (2);
\draw[violet,thick,-{Stealth}] (2,0) arc (0:50:2);
\draw[violetdeep,thick,dash pattern=on 3pt off 2pt] (0,1) circle (1);
\fill[violet] (1,0) circle (2.2pt);
\fill[violet] (0,1) circle (2.2pt);
\fill[violet] (-1,0) circle (2.2pt);
\fill[violet] (0,-1) circle (2.2pt);
\node[below right] at (1,0) {$\small 1$};
\node[left] at (0,1) {$\small i$};
\node[below left] at (-1,0) {$\small-1$};
\node[right] at (0,-1) {$\small -i$};
\node[below] at (2,-0.03) {$\small 2$};
\node[anchor=west,violet] at (2.02,0.95) {$\small |z|=2$};
\end{tikzpicture}
\end{center}

\textbf{图 2 读法}：灰色实线圆是 $|z|=1$，四个实心点 $1,i,-1,-i$ 正好落在它上面；浅紫色实线圆是本题的围道 $|z|=2$（图右标出），四个点都在它内部；深紫色点线圆是圆心在 $i$、半径 $1$ 的圆，后面变式 1 用它作围道，它只圈住 $i$。图上一个点恰好落在某条圆上，就说明那条围道不能直接用留数定理。


\textbf{第四步 · 算留数}。用规则 II（$P(z)=z$，$Q'(z)=4z^{3}$）：
\[ \operatorname{Res}f(z_k)=\frac{z_k}{4z_k^{3}}=\frac{1}{4z_k^{2}} . \]
代入：$z=\pm1$ 时 $z^{2}=1$，留数 $\frac14$；$z=\pm i$ 时 $z^{2}=-1$，留数 $-\frac14$。四个值是 $\frac14,\ \frac14,\ -\frac14,\ -\frac14$，和为 $0$。

\textbf{第五步 · 用留数定理}。
\[ I=2\pi i\sum_k\operatorname{Res}f(z_k)=2\pi i\cdot0=0 . \]

\aside{留数定理的完整形式：若 $f$ 在简单闭曲线 $C$ 内部除有限个孤立奇点 $z_1,\dots,z_n$ 外解析、且连续到 $C$ 上，则 $\oint_Cf\dd z=2\pi i\sum_{k=1}^{n}\operatorname{Res}f(z_k)$，其中 $C$ 取正向（逆时针）。本例的 $f\sim z^{-3}$，四个奇点都在内部，$C$ 上无奇点，条件满足。}

\textbf{复核一 · 无穷远留数}。$f$ 在无穷远像 $z^{-3}$，故 $\operatorname{Res}[f,\infty]=0$；而全平面留数和满足 $\sum_k\operatorname{Res}f(z_k)+\operatorname{Res}[f,\infty]=0$，所以四个有限留数之和必须为 $0$，与上面一致。

\textbf{复核二 · 数值}。沿 $|z|=2$ 直接做数值积分，结果的模小于 $10^{-14}$，说明它是 $0$，不是“接近 $0$ 的某个小数”。

\needspace{5\baselineskip}
\textbf{三个变式}（同一函数，只改围道）：

\textbf{变式 1}，$|z-i|=1$：圆心在 $i$、半径 $1$。$i$ 在内部；$1$ 与 $-1$ 到 $i$ 的距离都是 $\sqrt2>1$，$-i$ 的距离是 $2$，都在外部。只有一个极点，于是
\[ I=2\pi i\left(-\frac14\right)=-\frac{\pi i}{2} . \]

\textbf{变式 2}，$|z|=1$：四个极点全在围道上，留数定理不能用（§7.5）。

\textbf{变式 3}，$|z|=3$：四个全在内部，仍得 $0$——半径再大也不会变，因为留数和本身就是 $0$。



\hdr{§9}{四道配套例题（各练一段）}

例 A 取自课件 p20；例 B、C、D 是按同一套方法自编的补充题（不是课件原题），分别练“分母有重因子”“四个根同模”“只有一部分根在圈内”三种情形。

\qbreak
\qstem{例 A（课件 p20 例 3）求 $f(z)=\dfrac{1}{(z^{2}+1)^{2}}$ 在其奇点处的留数。}

\textbf{极点在哪个、几阶}：$z^{2}+1=0$ 给 $\pm i$；分母是平方，$(z^{2}+1)^{2}=(z-i)^{2}(z+i)^{2}$，所以两个都是二阶极点（§7.1 中 $k=0$ 那一支）。

\textbf{留数}（$m=2$，规则 III）：
\[ \operatorname{Res}f(i)=\lim_{z\to i}\frac{\dd}{\dd z}\left[(z-i)^{2}f(z)\right]
=\lim_{z\to i}\frac{\dd}{\dd z}\frac{1}{(z+i)^{2}}
=\frac{-2}{(2i)^{3}}=-\frac{i}{4}, \]
同法 $\operatorname{Res}f(-i)=\dfrac{i}{4}$。

\textbf{复核}：两留数之和为 $0$，与 $f\sim z^{-4}$（无穷远留数为 $0$）一致。

\qbreak
\qstem{例 B 设 $f(z)=\dfrac{1}{z^{3}-3z+2}$，写出全部极点与阶数，并求 $\displaystyle\oint_{|z|=2.5}f\dd z$。}

\textbf{极点}：$z^{3}-3z+2=(z-1)^{2}(z+2)$（§5.4），所以 $z=1$ 是二阶极点，$z=-2$ 是一阶极点。

\textbf{留数}：$\operatorname{Res}f(-2)=\dfrac{1}{Q'(-2)}=\dfrac{1}{12-3}=\dfrac19$；
\[ \operatorname{Res}f(1)=\lim_{z\to1}\frac{\dd}{\dd z}\frac{1}{z+2}=-\frac{1}{(1+2)^{2}}=-\frac19 . \]

\textbf{积分}：两个极点都在 $|z|=2.5$ 内（$|1|=1$，$|-2|=2$），留数之和为 $0$，故
\[ \oint_{|z|=2.5}f\dd z=2\pi i\cdot0=0 . \]

\textbf{复核}：$f\sim z^{-3}$，无穷远留数为 $0$，所以有限留数之和必为 $0$——这一条可以省掉逐个算。

\qbreak
\qstem{例 C 设 $f(z)=\dfrac{1}{z^{4}+4}$，写出四个极点，并求 $\displaystyle\oint_{|z|=2.5}f\dd z$。}

\textbf{极点}：$z^{4}=-4$。$a=-4=4e^{i\pi}$，模 $4^{1/4}=\sqrt2$，角 $\frac{\pi+2k\pi}{4}$，四个根是
\[ \sqrt2\,e^{i\pi/4}=1+i,\quad -1+i,\quad -1-i,\quad 1-i . \]
（与 §6.4 的 Sophie Germain 分解给出的根相同。）

\textbf{判圈}：四个根的模都是 $\sqrt2\approx1.414<2.5$，全部在内部。

\textbf{积分}：$f\sim z^{-4}$，无穷远留数为 $0$，故四个有限留数之和为 $0$，于是 $I=0$。若愿意逐个验算：四个留数依次为 $-\frac{1+i}{16},\ -\frac{1-i}{16},\ \frac{1-i}{16},\ \frac{1+i}{16}$，相加确实为 $0$。

\qbreak
\qstem{例 D 设 $f(z)=\dfrac{z}{z^{4}-5z^{2}+4}$，写出四个极点，并求 $\displaystyle\oint_{|z|=1.5}f\dd z$。}

\textbf{极点}：$z^{4}-5z^{2}+4=(z^{2}-1)(z^{2}-4)=(z-1)(z+1)(z-2)(z+2)$，四个一阶极点。

\textbf{判圈}：模分别是 $1,1,2,2$。围道 $|z|=1.5$ 只圈住 $\pm1$，$\pm2$ 在外部，不算。

\textbf{留数}（规则 II，$Q'(z)=4z^{3}-10z$）：
\[ \operatorname{Res}f(1)=\frac{1}{Q'(1)}=\frac{1}{-6}=-\frac16,
\qquad \operatorname{Res}f(-1)=\frac{-1}{Q'(-1)}=\frac{-1}{6}=-\frac16 . \]

\textbf{积分}：
\[ I=2\pi i\left(-\frac16-\frac16\right)=-\frac{2\pi i}{3} . \]

\textbf{复核}：如果把 $\pm2$ 误算进去，会再加上 $2\pi i\left(\frac16+\frac16\right)=+\frac{2\pi i}{3}$，正好把答案抵消成 $0$——这就是“圈外极点不能算”的具体代价。

\hdr{§10}{操作卡：四步}

拿到 $f=P/Q$ 与围道 $|z|=R$，按这四步走。

\begin{center}
\begin{tikzpicture}[node distance=6mm,
  box/.style={draw=violet,rounded corners=2pt,fill=bandviolet,text width=10.4cm,
              align=left,inner sep=5pt,font=\small}]
\node[box] (s1) {\textbf{步骤 1 · 分解分母}：$z^{n}=a$ 用式 \eqref{eq:roots}；结构型用 §6；都不行就试根降次（§5），一直分到一次或二次因子。};
\node[box,below=of s1] (s2) {\textbf{步骤 2 · 每个根记两件事}：模 $|z_k|$ 与重数 $m_k$；重数用导数检验（§1.2）或读因式指数（§7.2）。};
\node[box,below=of s2] (s3) {\textbf{步骤 3 · 与 $R$ 比}：小于留下，大于划掉，等于单独处理（§7.5）。};
\node[box,below=of s3] (s4) {\textbf{步骤 4 · 定阶算留数}：单根用规则 II；$m\ge2$ 用规则 III；分子分母同零点先约分（§7.1）。};
\draw[-{Stealth},violet] (s1) -- (s2);
\draw[-{Stealth},violet] (s2) -- (s3);
\draw[-{Stealth},violet] (s3) -- (s4);
\end{tikzpicture}
\end{center}

图 3 是这四步的流向。第 3 步如果出现“等于”，不要往下走，先转到 §7.5 决定是改围道还是挖半圆。

\hdr{§11}{练习（先做，详解见附录 A）}

六道题都是按本章方法自编的（不是课件作业），题干与考点都在这里，详解统一放在附录 A，方便先做后对。

\qbreak
\qstem{练习 1 解方程 $z^{4}=-16$，写出四个根的模与辐角；再判断这四个根相对 $|z|=1$、$|z|=2$、$|z|=3$ 的位置。}

\textbf{考点}：$z^{n}=a$ 的公式与“同模”判据。\label{ex:1}详解见附录 A 第 \pageref{ans:1} 页。

\qbreak
\qstem{练习 2 分解 $z^{4}-5z^{2}+4$，指出 $f(z)=\dfrac{z}{z^{4}-5z^{2}+4}$ 的四个极点各是几阶；再求 $\displaystyle\oint_{|z|=1.5}f\dd z$。}

\textbf{考点}：结构识别 + 只算圈内极点。\label{ex:2}详解见附录 A 第 \pageref{ans:2} 页。

\qbreak
\qstem{练习 3 求 $f(z)=\dfrac{1}{(z^{2}+1)^{2}}$ 的极点、阶数与留数。}

\textbf{考点}：分母是平方时的阶数判定与规则 III。\label{ex:3}详解见附录 A 第 \pageref{ans:3} 页。

\qbreak
\qstem{练习 4 判断 $z=1$ 是 $f(z)=\dfrac{z^{2}-1}{z^{4}-1}$ 的什么奇点，并说明理由。}

\textbf{考点}：分子分母同零点的约分（§7.1）。\label{ex:4}详解见附录 A 第 \pageref{ans:4} 页。

\qbreak
\qstem{练习 5 求 $f(z)=\dfrac{1}{z^{3}-3z+2}$ 的极点与阶数，并求 $\displaystyle\oint_{|z|=2.5}f\dd z$。}

\textbf{考点}：重根的识别与规则 II、III 的混用。\label{ex:5}详解见附录 A 第 \pageref{ans:5} 页。

\qbreak
\qstem{练习 6 例 2 的积分若改在 $|z|=1$ 上，还能直接用留数定理吗？说明原因，并给出两种处置办法。}

\textbf{考点}：根落在围道上的情形（§7.5）。\label{ex:6}详解见附录 A 第 \pageref{ans:6} 页。

\hdr{§12}{易错点（每条指回具体位置）}

\gapitem{一}{把 $z^{n}=-|a|$ 的辐角写成 $0$。$z^{4}=-1$ 若取 $\theta=0$，得 $\pm1,\pm i$，那是 $z^{4}=1$ 的根。见 §4.2。}

\gapitem{二}{试出一个根就停手。$z^{3}-3z+2$ 在 $z=1$ 处是二重根，只除一次会把二阶极点当一阶。见 §5.4。}

\gapitem{三}{把“分母是平方”直接当成“二阶极点”，不看分子是否同零点。$f=\dfrac{z^{2}-1}{z^{4}-1}$ 在 $z=1$ 处分子分母都为零，约分后是 $\dfrac{1}{z^{2}+1}$，$z=1$ 根本不是极点。见 §7.1。}

\gapitem{四}{开错次方，判错圈。$z^{4}=-16$ 的四根模是 $2$，不是 $16$，也不是 $\sqrt2$；要开四次方，不是开平方。见 §4.3 第四行。}

\gapitem{五}{把圈外的极点也算进去。$|z|=1.5$ 时 $\pm2$ 在圈外，算进去会把例 D 的答案抵消成 $0$。见例 D 的复核段。}

\gapitem{六}{根落在围道上还硬套留数定理。见 §7.5、练习 6。}

\hdr{§13}{缺口清单：课件跳了哪一步，本讲义补在哪}

\gapitem{代数缺口 · p19}{原文“函数 $f(z)=\frac{z}{z^{4}-1}$ 有四个一阶极点 $z=\pm1,\ \pm i$，且在积分半径内，则规则 II”直接给结果，没有解 $z^{4}-1=0$ 的过程。补法：§4 给出 $z^{n}=a$ 的通式与推导，§5.3 用试根法把同一组根走一遍。}

\gapitem{重数缺口 · p20}{原文“$z=\pm i$ 是该函数的两个二阶极点，因此有”，没有说明 $(z^{2}+1)^{2}=(z-i)^{2}(z+i)^{2}$ 的指数就是阶数。补法：§7.1 的 $m-k$ 阶判据，§7.2 的导数检验。}

\gapitem{压缩缺口 · p21}{原文“可以看出，该函数有两个极点，分别为三阶极点 $z=0$ 和一阶极点 $z=i$”，“可以看出”吞掉了分母各因式幂次的读法。补法：§7.2 检验一（因式分解读指数）。}

\gapitem{判圈缺口 · p22--p23}{原文只用两行不等式给出 $|z_1|>1>|z_2|$，没有交代依据。补法：§2.2 用韦达定理与“两根为负”补出这两行，§7.4 给出四种判圈情形。}

\gapitem{结构缺口 · p51}{原文把 $z^{4}+a^{4}=0$ 写成 $z^{4}=-a^{4}=a^{4}e^{i(2k+1)\pi}$ 并列出四个根，这一处写得完整，但没有指出它就是 $z^{n}=a$ 的通式。补法：§4 把式 \eqref{eq:roots} 作为统一公式给出，p51 成为它的一个实例。}

\gapitem{边界缺口 · p19、p22--p23}{两处都以 $|z|=2$、$|z|=1$ 一类围道作默认，没有说明“根恰在围道上该怎么办”。补法：§7.5、练习 6。}

\hdr{§14}{三句话}

第一句：\textbf{$n$ 次方程一共 $n$ 个根（按重数计）}，解完要数一遍；数目不对，就是漏了根或漏了重数。

第二句：\textbf{$z^{n}=a$ 的根全在半径 $|a|^{1/n}$ 的圆上、等角分布}，判圈一次到位；不是这个形状的，就识别结构或试根降次。

第三句：\textbf{极点的阶数等于分母零点的阶数减分子零点的阶数}，单根配规则 II，重根配规则 III，根在围道上先挖洞。

\hdr{§15}{与课件的页码对照}

\begin{minipage}{\linewidth}
\small
\tabtop
\begin{tabular}{@{}p{2.55cm}p{0.66\linewidth}@{}}
\thead{课件页} & \thead{本讲义位置} \\
\end{tabular}
\tabmid
\begin{tabular}{@{}p{2.55cm}p{0.66\linewidth}@{}}
p11--p12 & §7.3（规则 II 的完整陈述） \\
p18--p19 & §8（例 2 全程与三个变式） \\
p20 & §6.3、例 A（分母是平方 → 二阶极点） \\
p21 & §7.2（读出分母因式的指数） \\
p22--p23 & §2.2、§7.4（韦达定理判圈） \\
p49--p51 & §4.3 第四行、§6.4、例 C（$z^{4}=a$ 的四个根） \\
p63--p70 & §7.5（根在围道上，半留数） \\
\end{tabular}
\tabbot
\end{minipage}

\flowerrule

\hdr{附录 A}{练习详解}

每题先给考点与关键一步，再给详解、答案与自检；题号与正文 §11 一一对应。

\qbreak
\qstem{（正文练习 1，题面见第 \pageref{ex:1} 页）解 $z^{4}=-16$，写出四个根的模与辐角，并判断它们相对 $|z|=1$、$|z|=2$、$|z|=3$ 的位置。\label{ans:1}}

\textbf{考点}：$z^{n}=a$ 的公式（式 \eqref{eq:roots}）与“同模”判据。

\textbf{关键一步}：$a=-16$ 是负实数，$\theta=\pi$，不是 $0$。

\textbf{详解}：$|a|=16$，$4$ 次方根 $\sqrt[4]{16}=2$，所以所有根的模都是 $2$。辐角 $\frac{\pi+2k\pi}{4}$，即 $\frac{\pi}{4},\frac{3\pi}{4},\frac{5\pi}{4},\frac{7\pi}{4}$，四个根
\[ 2e^{i\pi/4}=\sqrt2(1+i),\quad \sqrt2(-1+i),\quad \sqrt2(-1-i),\quad \sqrt2(1-i). \]

\textbf{答案}：模全为 $2$；辐角 $45^\circ,135^\circ,225^\circ,315^\circ$。相对 $|z|=1$ 全在外部；相对 $|z|=2$ 四个都恰在围道上（不能用留数定理）；相对 $|z|=3$ 全在内部。

\textbf{自检}：四个根之和应为 $0$（$z^{3}$ 的系数为零），四个点关于原点对称，图上一看便知。

\qbreak
\qstem{（正文练习 2）分解 $z^{4}-5z^{2}+4$，指出 $f(z)=\dfrac{z}{z^{4}-5z^{2}+4}$ 各极点的阶数，并求 $\oint_{|z|=1.5}f\dd z$。\label{ans:2}}

\textbf{考点}：双二次分解（§6.2）与“只算圈内极点”（§7.4）。

\textbf{关键一步}：令 $u=z^{2}$，$u^{2}-5u+4=(u-1)(u-4)$。

\textbf{详解}：$z^{4}-5z^{2}+4=(z^{2}-1)(z^{2}-4)=(z-1)(z+1)(z-2)(z+2)$，四个互不相同的根，都是一阶极点。$Q'(z)=4z^{3}-10z$：$Q'(1)=-6$，$Q'(-1)=6$，都不为零，符合规则 II 的条件。圈内的只有 $\pm1$，留数各为 $-\frac16$，和为 $-\frac13$。

\textbf{答案}：$I=2\pi i\left(-\frac13\right)=-\dfrac{2\pi i}{3}$。

\textbf{自检}：把 $\pm2$ 误算进去会得到 $0$（见例 D 的复核段），所以“答案是不是 0”本身就是判圈是否正确的信号。

\qbreak
\qstem{（正文练习 3）求 $f(z)=\dfrac{1}{(z^{2}+1)^{2}}$ 的极点、阶数与留数。\label{ans:3}}

\textbf{考点}：分母为平方时的阶数（§7.1）与规则 III。

\textbf{关键一步}：$(z^{2}+1)^{2}=(z-i)^{2}(z+i)^{2}$，指数 $2$ 就是阶数。

\textbf{详解}：两个二阶极点 $\pm i$。$m=2$：
\[ \operatorname{Res}f(i)=\lim_{z\to i}\frac{\dd}{\dd z}\frac{1}{(z+i)^{2}}=\frac{-2}{(2i)^{3}}=\frac{-2}{-8i}=-\frac{i}{4}, \]
$\operatorname{Res}f(-i)=\dfrac{i}{4}$。

\textbf{答案}：$\pm i$ 都是二阶极点；留数分别为 $-\frac{i}{4}$ 与 $\frac{i}{4}$。

\textbf{自检}：两个留数之和为 $0$，与 $f\sim z^{-4}$ 一致；若误判成一阶极点去套规则 I，$\lim_{z\to i}(z-i)f$ 会发散到 $\infty$，立刻暴露阶数判错。

\qbreak
\qstem{（正文练习 4）判断 $z=1$ 是 $f(z)=\dfrac{z^{2}-1}{z^{4}-1}$ 的什么奇点，并说明理由。\label{ans:4}}

\textbf{考点}：分子分母同零点时的约分（式 \eqref{eq:order} 的 $k\ge m$ 支）。

\textbf{关键一步}：分子在 $z=1$ 是一阶零点，分母也是一阶零点，$k=m=1$。

\textbf{详解}：$z^{2}-1=(z-1)(z+1)$，$z^{4}-1=(z-1)(z+1)(z^{2}+1)$，约去 $(z-1)(z+1)$ 得
\[ f(z)=\frac{1}{z^{2}+1}\qquad(z\ne\pm1), \]
$z=1$ 处补定义 $f(1)=\frac12$ 后解析。

\textbf{答案}：$z=1$ 是\textbf{可去奇点}，不是极点；留数为 $0$。同理 $z=-1$ 也是可去奇点，$f$ 真正的极点是 $\pm i$。

\textbf{自检}：若只看到“分母为零”就判成一阶极点，用规则 II 会得到 $\frac{P(1)}{Q'(1)}=\frac{0}{4}=0$——这个 $0$ 其实是“分子也为零”的信号，不是留数。

\qbreak
\qstem{（正文练习 5）求 $f(z)=\dfrac{1}{z^{3}-3z+2}$ 的极点与阶数，并求 $\oint_{|z|=2.5}f\dd z$。\label{ans:5}}

\textbf{考点}：重根的识别（§5.4）与规则 II、III 的混用。

\textbf{关键一步}：$z^{3}-3z+2=(z-1)^{2}(z+2)$。

\textbf{详解}：$z=1$ 二阶、$z=-2$ 一阶。$\operatorname{Res}f(-2)=\frac{1}{Q'(-2)}=\frac19$；$\operatorname{Res}f(1)=\frac{\dd}{\dd z}\frac{1}{z+2}\big|_{1}=-\frac19$。两者都在 $|z|=2.5$ 内，和为 $0$。

\textbf{答案}：$I=0$。

\textbf{自检}：$f\sim z^{-3}$，无穷远留数为 $0$，故有限留数之和必为 $0$，与逐项计算一致。

\qbreak
\qstem{（正文练习 6）例 2 的积分若改在 $|z|=1$ 上，还能直接用留数定理吗？说明原因并给出两种处置办法。\label{ans:6}}

\textbf{考点}：根落在围道上的情形（§7.5）。

\textbf{关键一步}：$z^{4}=1$ 的四个根模都是 $1$，恰好全在 $|z|=1$ 上。

\textbf{详解}：留数定理要求奇点严格位于围道内部，这四个点都不满足，所以不能直接套 $2\pi i\sum\operatorname{Res}$。处置办法有两种：一是换围道，例如取 $|z|=R$（$R\ne1$），此时四个点全在内部或全在外部，积分分别为 $0$ 或按 $-2\pi i\sum$ 计算（本题两条都给 $0$）；二是在四个根处各挖一个小半圆，让围道从它们旁边绕过，取主值，此时每个一阶极点贡献 $i\pi\operatorname{Res}$，四项相加为 $i\pi\cdot0=0$。

\textbf{答案}：不能直接用；两种处置分别是换围道与挖小半圆取主值，结果都是 $0$。

\textbf{自检}：只要围道穿过极点，就必须先处理“穿点”再谈留数，否则答案可能差一个 $i\pi$ 项。

\flowerrule

\needspace{0.66\textheight}
\hdr{附录 B}{速查卡（可单独打印）}

\subhdr{记账与重数}

\begin{itemize}
\item $n$ 次方程一共 $n$ 个根（按重数计）；$(z^{2}+1)^{2}=0$ 是 $4$ 个根、$2$ 个不同值。
\item $P(z_0)=\cdots=P^{(m-1)}(z_0)=0$ 且 $P^{(m)}(z_0)\ne0$ $\Longrightarrow$ $m$ 重根；只看 $P'(z_0)\ne0$ 即可确认单根。
\item 实系数方程：非实根共轭成对、重数相同，解出一半即可。
\end{itemize}

\subhdr{求根三路}

\begin{itemize}
\item \textbf{二项式}：$z^{n}=a$ 的根为 $|a|^{1/n}e^{i(\theta+2k\pi)/n}$，$k=0,\dots,n-1$；同模、等角、正 $n$ 边形。
\item \textbf{结构}：平方差、立方差、双二次（令 $u=z^{2}$）、完全平方（重根就从这里来）、
$z^{4}+4=(z^{2}-2z+2)(z^{2}+2z+2)$。
\item \textbf{猜根}：整系数首一时，试常数项的因数；除 $(z-z_1)$ 降次；商在同一点仍为 $0$ 就再除一次。
\end{itemize}

\subhdr{四个必背的值}

\begin{itemize}
\item $z^{4}=1$：$1,\ i,\ -1,\ -i$。\quad $z^{4}=-1$：$\dfrac{1+i}{\sqrt2},\ \dfrac{-1+i}{\sqrt2},\ \dfrac{-1-i}{\sqrt2},\ \dfrac{1-i}{\sqrt2}$。
\item $z^{3}=8$：$2,\ -1+i\sqrt3,\ -1-i\sqrt3$。\quad $z^{4}=-16$：$\sqrt2(\pm1\pm i)$，模为 $2$。
\end{itemize}

\subhdr{从根到极点}

\begin{itemize}
\item 阶数 $=m-k$：分母零点阶 $m$ 减分子零点阶 $k$；$k\ge m$ 时可去，不是极点。
\item 单根用 $\operatorname{Res}=\dfrac{P(z_0)}{Q'(z_0)}$；$m\ge2$ 用规则 III 的导数式；同零点先约分。
\end{itemize}

\subhdr{判圈}

\begin{itemize}
\item 算 $|z_k|$ 与 $R$ 比；等于 $R$ 说明根在围道上，要先挖洞或换围道。
\item $z^{n}=a$ 一次判全体；共轭对用 $|z|^{2}=c/a$；实根用韦达的乘积与和。
\end{itemize}

\subhdr{四步（照做）}

分解分母 $\to$ 记每个根的模与重数 $\to$ 与 $R$ 比 $\to$ 定阶并算留数。

\subhdr{自检三问}

\begin{enumerate}
\item 根的数量对得上次数吗（按重数计）？
\item 每个根的模都比过 $R$ 了吗？有没有恰好落在围道上的？
\item 留数用的规则对吗（可去、单根、重根三种情形各不同）？
\end{enumerate}

\end{document}
